\documentclass[review,12pt]{elsarticle}
\usepackage{amsmath,amssymb,amsfonts}
\usepackage{graphicx}
\usepackage{booktabs}
\usepackage{tabularx}
\usepackage{xcolor}
\usepackage[hidelinks]{hyperref}
\usepackage{xurl}
\usepackage{xr}
\externaldocument{main}

\journal{Knowledge-Based Systems}

\begin{document}
\begin{frontmatter}
\title{Supplementary Material for DH$^2$A-KT: Leakage-Audited
Knowledge Attribution for Knowledge Tracing}
\author{Tuan Dao Minh and co-authors}
\end{frontmatter}

These notes are exploratory. They are not research questions and not
contributions of the main article. Compile this file after
\texttt{main.tex} so cross-references resolve.

\section{Exploratory Critic-gated explanations}
\label{sec:tier2-supp}
Three agents read frozen DH$^2$-KT scores: a Diagnostician writes a
rationale, a Critic flags output that contradicts the frozen
$P(\mathrm{correct})$, and a Tutor writes session hyperedges that the
frozen tracer does not consume. We ran $n{=}500$ prompts (XES3G5M
fold~0, seed 42, Qwen2.5-7B via Ollama) on the graph-only v2 checkpoint
and again on QKC-T. Independent generation (IG) omits
$P(\mathrm{correct})$ from the Diagnostician prompt while the Critic
still sees the true score, so the flag-rate gap is largely definitional
and is not a content finding.
On QKC-T, grounded/IG flags are $0.018$/$0.766$ with mean JC
$0.161$/$0.152$; the JC gap is $\approx{+}0.010$.
Blinded Likert ratings on a seed~42 subsample of $n{=}100$ grounded
tracer explanations (3 independent raters) give faithfulness
$4.55{\pm}0.63$ and usefulness $4.01{\pm}0.74$ (ordinal Krippendorff
$\alpha{=}0.52$/$0.56$; mean pairwise Spearman $\rho{=}0.43$/$0.34$;
within$\pm$1 $95.0\%$/$94.7\%$). An earlier dual-rater $n{=}40$ pack on
the v2 graph-only checkpoint is archived and is not the primary
human-eval row.
Removing the Critic sets the flag rate to $0$ by construction while
leaving ID-anchored KC-Jaccard unchanged.
% Auto-derived from results/tables/*tier2_pilot*_summary.json
\begin{table}[!t]
\caption{Tier~2 agent faithfulness pilot (XES3G5M fold~0, $n{=}500$, seed 42, Qwen2.5-7B via Ollama). Flag rate = fraction of Diagnostician outputs flagged by the mandatory Critic. v2 applies probability-format calibration (\texttt{critic\_calibration.py}).}
\label{tab:tier2}
\centering
\begin{tabular}{@{}lccc@{}}
\toprule
Run & Backend & Flag rate & Mean $P(\mathrm{correct})$ \\
\midrule
Stub (sanity) & StubLLM & 0.000 (0/500) & 0.742 \\
Ollama v1 & qwen2.5:7b & 0.124 (62/500) & 0.746 \\
Ollama v2 & qwen2.5:7b & \textbf{0.000} (0/500) & 0.753 \\
\bottomrule
\end{tabular}
\end{table}

% Auto-updated by scripts/19_update_kbs_faithfulness_table.py
\begin{table}[!t]
\caption{Tier~2 explanation faithfulness: grounded Diagnostician vs.\ independent
generation (IG). ID-anchored KC-Jaccard (JC) measures overlap between
\texttt{GroundedKCs} trailer IDs and the support set
$\{\mathrm{target}\}\cup\mathrm{history}\cup 1$-hop $E_{\mathrm{pre}}$.
IG omits frozen $P(\mathrm{correct})$ from the Diagnostician prompt; the Critic
still sees the true Tier-1 score.}
\label{tab:kc-jaccard-ig}
\centering
\begin{tabular}{@{}llccc@{}}
\toprule
Arm & Backend & $n$ & Flag rate & Mean JC \\
\midrule
Grounded & StubLLM (harness) & 20 & 0.000 & 0.191 \\
IG & StubLLM (harness) & 20 & 1.000 & 0.191 \\
Grounded & qwen2.5:7b (Ollama) & 500 & 0.026 & 0.152 \\
IG & qwen2.5:7b (Ollama) & 500 & 0.768 & 0.142 \\
\bottomrule
\end{tabular}
\end{table}

\begin{table}[!t]
\caption{Human rating of Tier~2 outputs on the recommended tracer QKC-T
(XES3G5M fold~0, grounded Ollama subsample, $n{=}100$ explanations,
3 independent raters, Likert 1--5). Agreement is the mean of pairwise
Spearman $\rho$ and the mean fraction of scores within $\pm 1$.}
\label{tab:human-eval}
\centering
\begin{tabular}{@{}lcc@{}}
\toprule
Scale & Mean $\pm$ SD & Agreement (3 raters, mean pairwise) \\
\midrule
Faithfulness & $4.55 \pm 0.63$ & $\rho{=}0.43$; within$\pm$1: 95.0\% \\
Usefulness & $4.01 \pm 0.74$ & $\rho{=}0.34$; within$\pm$1: 94.7\% \\
\bottomrule
\end{tabular}
\end{table}

% Paired ablation: predictive (Tier 1) + faithfulness (Tier 2) knobs.
% Numbers from results/tables/* (Ollama n=500; FA session ablation; manip folds).
\begin{table}[!t]
\caption{Paired ablation of DH$^2$A-KT components. Tier~1 metrics are on
XES3G5M unless noted; Tier~2 faithfulness uses the fold~0 frozen checkpoint
and Qwen2.5-7B (Ollama, $n{=}500$) unless noted. ``w/o Critic'' disables the
gate on the same grounded Diagnostician outputs (flag rate $0$ by construction;
JC unchanged). Independent generation (IG) strips frozen $P(\mathrm{correct})$
from the Diagnostician. Session hyperedges are FoundationalASSIST-only.
``w/o $L_{\mathrm{aux}}$'' sets \texttt{graph\_sensitivity\_weight}${}=0$
(fold~0 retrain).}
\label{tab:paired-ablation}
\centering
\footnotesize
\begin{tabular}{@{}llccl@{}}
\toprule
Setting & AUC & Manip.\ $\Delta$AUC & Flag & Mean JC \\
\midrule
Full (grounded + Critic) & $0.752{\pm}0.002$ & $0.038$--$0.047$ & $0.026$ & $0.152$ \\
IG (w/o $P$ in Diag.) & --- & --- & $0.768$ & $0.142$ \\
w/o Critic & --- & --- & $0.000$ & $0.152$ \\
FA: concept-only & $0.720$ & --- & --- & --- \\
FA: + session HE & $0.722$ & --- & --- & --- \\
w/o $L_{\mathrm{aux}}$ (fold~0) & $0.754$ & $0.039$ & --- & --- \\
\bottomrule
\end{tabular}
\end{table}


\section{Observational hint identification check}
\label{sec:ate-supp}
On FoundationalASSIST we report a standard inverse-propensity (IPW)
average treatment effect of $\texttt{hint\_used}_t$ on
$\texttt{correct}_{t+1}$, with logistic propensity on pre-$t$ summaries
and clipping. Unconfoundedness is assumed, not proven. Same-step
\texttt{discrete\_score} is not the outcome (it is mechanically tied to
hint requests). No estimate is reported on XES3G5M (no hint logs).

Using train transitions, the IPW estimate of hint use on
\emph{next-step} correctness is $-0.207$ (naive mean difference
$-0.255$; trimmed IPW $-0.211$; bootstrap 95\% CI
$[-0.230,\,-0.187]$ on 50k-row resamples;
$n_{\mathrm{treated}}{=}36{,}151$). Propensity scores hit the clip floor
($[0.05,\,0.219]$), indicating weak overlap. The negative sign is
\emph{consistent with} selection and confounding (weaker learners
request more hints); it is not identified as the effect of providing
hints, and it is not a recommendation against hints. We did not compute
Rosenbaum bounds or an E-value.
\begin{table}[htbp]
\caption{FoundationalASSIST fold~0 IPW identification check of
\texttt{hint\_used}$_t$ on \texttt{correct}$_{t+1}$ under pre-$t$
confounders. Propensity mass hits the clip floor. The negative IPW
estimate reflects that weaker learners request more hints; it is
\emph{not} the effect of providing hints.}
\label{tab:hint-ate}
\centering
\begin{tabular}{@{}lr@{}}
\toprule
Quantity & Value \\
\midrule
Hint ATE (IPW) / naive mean diff & $-0.207$ / $-0.255$ \\
\quad trimmed IPW (propensity 10--90\%) & $-0.211$ \\
\quad bootstrap 95\% CI (subsample) & $[-0.230,\,-0.187]$ \\
\quad $n$ treated / control (train transitions) & 36{,}151 / 660{,}321 \\
\quad propensity $[\min,\max]$ & $[0.05,\,0.219]$ \\
\bottomrule
\end{tabular}
\end{table}


\end{document}
